\chapter{Force, mobility, flux and finite action value}\label{ch:force}
The potential describes a state. Its derivative describes how that potential
changes in a direction. A constitutive response law describes how strongly
a controller moves in response. A finite EBU quote describes the complete
change caused by a specified action. These four statements fit together,
but none should be substituted for another.

\section{The force along a physical edge}
For a lossless edge $i\to j$, the unit action column is
$S_e=-e_i+e_j$, where $e_i$ is the coordinate vector with one at position
$i$. The directional derivative is
\begin{equation}
 \mu^T S_e=-\mu_i+\mu_j.
\end{equation}
Define the field force with the sign that makes a small improvement positive:
\begin{equation}\label{eq:force}
 f_e=-S_e^T\mu=\mu_i-\mu_j.
\end{equation}
For the Gaussian potential this is
\begin{equation}
 f_e=\frac{x_i-x_i^*}{\sigma_i^2}
       -\frac{x_j-x_j^*}{\sigma_j^2}.
\end{equation}
It compares standardized marginal deviations, not stocks alone. Two cells
holding equal quantities can have different marginals if their references
or scales differ. Conversely, unequal stocks can have equal marginals.

In the historical loss-aware edge model $S_e=-e_i+\eta_e e_j$, the same
chain rule gives $f_e=\mu_i-\eta_e\mu_j$. The efficiency belongs inside
the derivative because it changes the physical transformation. This general
directional identity does not by itself adopt a loss-aware runtime or
resolve its separate conservation and cost declarations.

\section{Mobility and the Onsager-type response}
The historical reference is Onsager's 1931 work on coupled irreversible
processes \cite{onsager-one,onsager-two}. It connects linear force--rate
relations and reciprocal coefficients under microscopic-reversibility
assumptions. EBU borrows a modelling structure, not an entitlement to those
physical assumptions. A chosen deviation potential need not be free energy,
and a declared mobility need not be an experimentally measured transport
coefficient. This is why the qualifier ``Onsager-type'' matters.

The reason for introducing the structure is constructive: it separates the
state's driving slope from a mechanism's responsiveness, and it permits an
explicit conditional dissipation calculation. It does not follow from the
finite accounting identity. We can change the response law while retaining
that identity, but we cannot then retain the old response law's dynamical proof.

The original EBU books studied the directed threshold-mobility law
\begin{equation}\label{eq:onsager}
 J_e=M_e[f_e-\theta_e]_+,\qquad M_e>0,\quad\theta_e\geq0,
\end{equation}
where $[a]_+=\max(a,0)$. $J_e$ is a rate. $M_e$ is mobility: the strength
of the declared response to force above threshold. In the present dimensionless
potential normalization, $M_e$ has units $\X^2/T$ and $\theta_e$ has units
$1/\X$. They are not the physical stock capacity or the persistent EBU
capacity balance introduced later.

This is the project's \emph{Onsager-type} constitutive controller. The
unthresholded signed relation $J=Mf$ is linear; the positive-part directed
version is piecewise linear and blocks reverse flow on that edge. Neither
form is derived merely by defining $V$. Choosing it is a separate modelling
decision. The name does not establish universal thermodynamic validity for
an economic system, nor does it imply a general coupled reciprocal-mobility
matrix. The historical proofs concern their explicitly stated diagonal,
directed controller.

If $M=2$, $f=1.4$ and $\theta=0.5$, the requested rate is $1.8$ in the
corresponding units. Doubling mobility doubles this unconstrained request.
It does not double the stock available at the source or authorize a physical
cap to be exceeded. Permission may produce a different accepted rate, in
which case a proof about the unconstrained law needs reinspection.

\section{The exact finite quadratic}
Let a declared action have finite increment $\delta$. Expanding the square,
\begin{align}
 V(z+\delta)
 &=\tfrac12(z-x^*+\delta)^TH(z-x^*+\delta)\\
 &=V(z)+\mu(z)^T\delta+\tfrac12\delta^TH\delta.
\end{align}
The two cross terms are equal because $H$ is symmetric. Therefore
\begin{equation}\label{eq:finite}
 E(\delta)=V(z)-V(z+\delta)-C
 =-\mu(z)^T\delta-\tfrac12\delta^TH\delta-C.
\end{equation}
No small-action approximation was made. The first term is the tangent
benefit; the second is the exact curvature correction. For $\delta=qS_e$
on a lossless edge,
\begin{equation}\label{eq:edgefinite}
 E(q)=qf_e(z)-\frac{q^2}{2}
       \left(\sigma_i^{-2}+\sigma_j^{-2}\right)-C(q).
\end{equation}
If $C=0$ and $f_e>0$, positive quantities give nonnegative value only up
to $2f_e/(\sigma_i^{-2}+\sigma_j^{-2})$. Physical constraints can be
tighter. The mathematical maximizer is half that upper zero, but the active
random-affordable programme does not automatically choose the maximizer.

\section{A force can be positive while a large action is harmful}
At $z=(18,2)$, $x^*=(10,10)$ and $\sigma=(2,2)$, the marginals are
$(2,-2)$, so $f=4$ and
\begin{equation}
 E(q)=4q-q^2/4.
\end{equation}
At $q=4$ the value is twelve. At $q=8$ it is sixteen. At $q=16$ it is
zero, and at $q=17$ it is $-17/4$. The initial force stays four in the
frozen tangent calculation, but the field changes along the finite action.
Quoting $qf$ would ignore that change and overstate benefit.

The Onsager controller would supply a rate based on the force and its own
mobility/threshold. A duration turns that rate into a quantity. Exact EBU
then values the resulting accepted quantity. Keeping these steps explicit
prevents the words force, flux and value from collapsing into one number.

\section*{Chapter recap}
Force is a directional derivative; mobility is a constitutive response
parameter; flux is the controller's rate; finite EBU is an endpoint
difference. Chapter~\ref{ch:dissipation} derives the dissipation result for the Onsager-type
law and explains exactly why it is not a recovery proof for a different actor.
